\documentclass[12pt, a4paper]{article}

\usepackage{amsmath, amssymb}
\usepackage{fontspec}
\setmainfont{FreeSerif}
\usepackage{geometry}
\geometry{margin=1in}
\usepackage[colorlinks=true, linkcolor=blue, urlcolor=blue, citecolor=blue]{hyperref}

\setcounter{secnumdepth}{0} 

\begin{document}

\section{Հավանականությունների տեսություն (06.02.26)}

\tableofcontents
\vspace{2em}
\hrule
\vspace{2em}

\subsection{\textbf{I. Լրիվ հավանականության բանաձև}}

Դիցուք ունենք $B$ պատահույթը և $A_1, A_2, \dots, A_n$ պատահույթները, որոնք կազմում են \textbf{լրիվ խումբ}, այսինքն՝

\begin{enumerate}
    \item $A_i \cap A_j = \emptyset$ (երբ $i \neq j$)
    \item $B$ պատահույթը կարող է հանդես գալ $A_i$-երից միայն մեկի հետ։
\end{enumerate}

Այսինքն՝ $B = (A_1 \cap B) \cup (A_2 \cap B) \cup \dots \cup (A_n \cap B)$

Քանի որ $(A_i \cap B) \cap (A_j \cap B) = \emptyset$ (որովհետև $A_i \cap A_j = \emptyset$), ապա ըստ գումարման թեորեմի.

$$P(B) = P(A_1 \cap B) + P(A_2 \cap B) + \dots + P(A_n \cap B)$$

Օգտվելով բազմապատկման թեորեմից՝ $P(A_i \cap B) = P(A_i) \cdot P(B|A_i)$, ստանում ենք \textbf{լրիվ հավանականության բանաձևը}.

$$P(B) = \sum_{i=1}^{n} P(A_i) \cdot P(B|A_i)$$

\vspace{1em}
\hrule
\vspace{1em}

\subsection{\textbf{II. Բայեսի ընդհանրացված բանաձև}}

$A_1, A_2, \dots, A_n$ պատահույթները կոչվում են \textbf{հիպոթեզներ}:

Քանի որ $P(A_i \cap B) = P(A_i) \cdot P(B|A_i) = P(B) \cdot P(A_i|B)$, ապա.

$$P(A_i|B) = \frac{P(A_i \cap B)}{P(B)} = \frac{P(A_i) \cdot P(B|A_i)}{\sum_{i=1}^{n} P(A_i) \cdot P(B|A_i)}$$

Սա կոչվում է \textbf{Բայեսի ընդհանրացված բանաձև}:

\vspace{1em}
\hrule
\vspace{1em}

\subsection{\textbf{III. Հավանականությունների աքսիոմատիկ սահմանումը}}

Դիցուք $X$ բազմության ենթաբազմությունների $S$ ոչ դատարկ բազմությունը կոչվում է \textbf{բազմությունների հանրահաշիվ}, եթե բավարարում է հետևյալ պայմաններին.

\begin{itemize}
    \item $A, B \in S \implies A \cap B \in S$
    \item $A \in S \implies \bar{A} = X \setminus A \in S$
\end{itemize}

\textbf{Սահմանում:} Բազմությունների հանրահաշիվը կոչվում է \textbf{$\sigma$-հանրահաշիվ}, եթե այն փակ է հաշվելի թվով միավորում գործողության նկատմամբ.

$$A_1, A_2, \dots, A_n, \dots \in S \implies \bigcup_{i=1}^{\infty} A_i \in S$$

Ըստ Դե Մորգանի օրենքի՝,
$$\overline{\bigcup_{i=1}^{\infty} A_i} = \bigcap_{i=1}^{\infty} \bar{A_i}$$հետևաբար հատումը նույնպես պատկանում է $S$-ին։

\textbf{Սահմանում:} $S$ $\sigma$-հանրահաշիվը կոչվում է \textbf{միավորով օժտված}, եթե $\exists \Omega \in S$ այնպիսին, որ $\forall A \in S$ համար $A \cap \Omega = A$:

\textbf{Հավանականային տարածություն:} $(\Omega, F, P)$ եռյակը, որտեղ.

\begin{itemize}
    \item $\Omega$-ն փորձի էլեմենտար ելքերի բազմությունն է,
    \item $F$-ը $\Omega$-միավորով $\sigma$-հանրահաշիվն է,
    \item $P$-ն հավանականային ֆունկցիան է $P: F \to [0, 1]$:
\end{itemize}

\vspace{1em}
\hrule
\vspace{1em}

\subsection{\textbf{IV. Հավանականության հատկությունները}}

\begin{enumerate}
    \item $P(\emptyset) = 0, P(\Omega) = 1$
    \item $P(\bigcup_{i=1}^{\infty} A_i) = \sum_{i=1}^{\infty} P(A_i)$, որտեղ $A_i \cap A_j = \emptyset \quad$($\sigma$-ադիտիվություն)
    \item $P(A \cup B) = P(A) + P(B) - P(A \cap B)$
\end{enumerate}

\subsubsection{\textbf{Բազմաչափ (Պոլինոմիալ) փորձերի սխեմա}}

Դիտարկենք $n$ անկախ փորձեր, որոնցից յուրաքանչյուրն ունի $k$ հնարավոր ելքեր՝ $A_1, A_2, \dots, A_k$ պատահույթների տեսքով:

\textbf{Հավանականություններ}. Յուրաքանչյուր $A_i$ պատահույթի հայտնվելու հավանականությունը հաստատուն է և հավասար է $p_i$, այսինքն՝ $P(A_i) = p_i$:

\textbf{Պայման.} Քանի որ սրանք կազմում են լրիվ խումբ, ապա $\sum_{i=1}^k p_i = 1$։

\textbf{Փորձերի արդյունքը.} Մեզ հետաքրքրում է այն դեպքը, երբ $n$ փորձերի արդյունքում $A_1$ պատահույթը հանդես է գալիս $m_1$ անգամ, $A_2$-ը՝ $m_2$ անգամ, ..., $A_k$-ն՝ $m_k$ անգամ:

Ընդհանուր քանակ. $m_1 + m_2 + \dots + m_k = n$:

Այն հավանականությունը, որ $n$ անկախ փորձերում $A_1, \dots, A_k$ պատահույթները հանդես կգան համապատասխանաբար $m_1, \dots, m_k$ անգամ, որոշվում է հետևյալ բանաձևով:

$$P_{n}(m_1, m_2, \dots, m_k) = \frac{n!}{m_1! m_2! \dots m_k!} \cdot p_1^{m_1} p_2^{m_2} \dots p_k^{m_k}$$

\textbf{Պոլինոմիալ գործակից.} $\frac{n!}{m_1! m_2! \dots m_k!}$ — սա ցույց է տալիս, թե քանի ձևով կարելի է $n$ տարրերը բաժանել $k$ խմբերի՝ համապատասխանաբար $m_1, \dots, m_k$ քանակներով:

\textbf{Հավանականային մաս}. $p_1^{m_1} p_2^{m_2} \dots p_k^{m_k}$ — սա կոնկրետ մեկ հաջորդականության հավանականությունն է, որտեղ յուրաքանչյուր ելք հանդես է գալիս իր քանակով:

\textbf{Կապը Բեռնուլիի բանաձևի հետ}
Եթե $k=2$ (այսինքն՝ ունենք միայն «հաջողություն» $A$ և «ձախողում» $\bar{A}$), ապա բանաձևը վերածվում է Բեռնուլիի բանաձևին.

$$P_n(m) = \frac{n!}{m!(n-m)!} p^m q^{n-m}$$

\textbf{Բեռնուլիի բանաձև:}

Եթե կատարվում է $n$ անկախ փորձ, և յուրաքանչյուրում $A$ պատահույթի հայտնվելու հավանականությունը $P(A)=p$, իսկ չհայտնվելունը՝ $P(\bar{A})=q=1-p$, ապա $m$ հաջողության հավանականությունը.

$$P_m(A) = C_n^m p^m q^{n-m}$$

\vspace{1em}
\hrule
\vspace{1em}

\subsection{\textbf{V. Հավանականության երկրաչափական մեկնաբանություն}}

\begin{itemize}
    \item Տարրական ելքերի տարածություն ($\Omega$): Օրինակ՝ $\Omega = [0, 1]$ հատվածը:
    \item $\sigma$-հանրահաշիվ ($F$): $\Omega$-ն միավորող $\sigma$-հանրահաշիվն է, որտեղ յուրաքանչյուր $A \in F$ պատահույթի համար $A \cap \Omega = A$:
    \item \textbf{Երկրաչափական հավանականություն:} $A$ պատահույթի հավանականությունը որոշվում է նրա չափի ($\mu$) հարաբերությամբ ամբողջ տիրույթի չափին.
\end{itemize}

$$P(A) = \frac{\mu(A)}{\mu(\Omega)}$$

\textbf{Կարևոր նկատառում:} Եթե բազմությունը չափելի չէ (օրինակ՝ Վիտալիի $E$ բազմությունը), ապա նրա չափը՝ $\mu(E)$, գոյություն չունի ($\nexists$), և հետևաբար հնարավոր չէ հաշվել դրա հավանականությունը:

\vspace{1em}
\hrule
\vspace{1em}

\subsection{\textbf{VI. Հիպերերկրաչափական բաշխում}}

Ունենք $N$ տարրից կազմված համախումբ, որից $M$-ը օժտված է $X$ հատկությամբ: Ընտրում ենք $n$ տարր: Հավանականությունը, որ դրանցից $m$-ը ($m \leq n$) կունենան $X$ հատկությունը.

$$P(A) = \frac{C_M^m \cdot C_{N-M}^{n-m}}{C_N^n}$$

\textit{Օրինակ:} 100 դետալ, 30-ը կանաչ են, 70-ը՝ ոչ: Ընտրում ենք 40-ը: Հավանականությունը, որ 40-ից 10-ը կանաչ են՝ $P(A) = \frac{C_{30}^{10} \cdot C_{70}^{30}}{C_{100}^{40}}$:

\vspace{1em}
\hrule
\vspace{1em}

\subsection{\textbf{VII. Պատահույթների երևումների ամենահավանական թիվը}}

Եթե մետաղադրամը նետում ենք $n=10$ անգամ, և մեզ հետաքրքրում է «զինանշան» գալու հավանականությունը $m$ անգամ, ապա քայլերը հետևյալն են.

\begin{itemize}
    \item Փորձերի քանակը ($n$): 10.
    \item Հաջողության հավանականությունը ($p$): Քանի որ մետաղադրամը սովորական է, զինանշան գալու հավանականությունը $1/2$ է.
    \item Ձախողման հավանականությունը ($q$): Նույնպես $1/2$ (քանի որ $q = 1 - p$).
\end{itemize}

\textbf{Հավանականության հաշվարկը ($m$ անգամի համար)}
Ըստ Բեռնուլիի բանաձևի, հավանականությունը, որ զինանշան կգա ուղիղ $m$ անգամ, հաշվվում է այսպես.

$$P_{10}(m) = C_{10}^m \cdot \left(\frac{1}{2}\right)^m \cdot \left(\frac{1}{2}\right)^{10-m}$$

Ուղիղ 1 անգամ ($m=1$):

$$P_{10}(1) = \frac{10!}{1!(10-1)!} \cdot \left(\frac{1}{2}\right)^1 \cdot \left(\frac{1}{2}\right)^9 = 10 \cdot \left(\frac{1}{2}\right)^{10}$$

Ուղիղ 2 անգամ ($m=2$):

$$P_{10}(2) = \frac{10!}{2!(10-2)!} \cdot \left(\frac{1}{2}\right)^2 \cdot \left(\frac{1}{2}\right)^8 = 45 \cdot \left(\frac{1}{2}\right)^{10}$$

\textbf{Ամենահավանական թիվը}
Հաջորդ տրամաբանական հարցն այն է, թե քանի՞ անգամ ամենայն հավանականությամբ կգա զինանշանը: Քանի որ $p=1/2$, ամենահավանական թիվը կլինի 5 անգամը: Դու կարող ես շարունակել հաշվարկը $P_{10}(5)$-ի համար և կտեսնես, որ այն ամենամեծ արժեքն է ստանում:

\textbf{Կապը բազմաչափ փորձերի հետ}
Եթե մետաղադրամի փոխարեն նետեիր զառ, որն ունի 6 տարբեր ելքեր, ապա պետք է օգտագործեիր քո նախորդ նկարների բազմաչափ (պոլինոմիալ) բանաձևը: Այդ դեպքում համարիչում կլիներ $n!$, իսկ հայտարարում՝ յուրաքանչյուր թվի հայտնվելու քանակների ֆակտորիալների արտադրյալը:

\vspace{1em}
\hrule
\vspace{1em}

\textbf{Խնդիր.} $n$ անկախ փորձերում, որոնցից յուրաքանչյուրում հաջողության հավանականությունը $p$ է, գտնել այնպիսի $m_0$ թիվ, որի դեպքում $P_n(m)$ հավանականությունն ընդունում է իր առավելագույն արժեքը:

Որպեսզի հասկանանք $P_n(m)$ ֆունկցիայի վարքը, դիտարկենք հաջորդական անդամների հարաբերությունը.

$$\frac{P_n(m+1)}{P_n(m)} = \frac{n! \cdot p^{m+1} \cdot q^{n-m-1}}{(m+1)!(n-m-1)!} \cdot \frac{m!(n-m)!}{n! \cdot p^m \cdot q^{n-m}} = \frac{p(n-m)}{q(m+1)}$$

\begin{itemize}
    \item Հավանականությունն աճում է, եթե $\frac{P_n(m+1)}{P_n(m)} > 1$: Սա տեղի ունի, երբ $m < np - q$:
    \item Հավանականությունը նվազում է, եթե $\frac{P_n(m+1)}{P_n(m)} < 1$: Սա տեղի ունի, երբ $m > np - q$:
\end{itemize}

\textbf{Ամենահավանական թվի ($m_0$) միջակայքը}
Այս երկու պայմանների համադրումից բխում է, որ $m_0$ թիվը պետք է բավարարի հետևյալ կրկնակի անհավասարմանը.

$$np - q \leq m_0 \leq np + p$$

\textbf{Հնարավոր դեպքերը}
Քանի որ միջակայքի երկարությունը հավասար է $1$-ի ($(np+p) - (np-q) = p+q = 1$), ապա հնարավոր է երկու իրավիճակ.

\begin{itemize}
    \item Եթե $np-q$ թիվը ամբողջ չէ. Միջակայքում կա ճիշտ մեկ ամբողջ թիվ: Դա էլ հանդիսանում է միակ ամենահավանական թիվը:
    \item Եթե $np-q$ թիվը ամբողջ է. Միջակայքի ծայրակետերը ($m_0 = np-q$ և $m_0+1 = np+p$) երկուսն էլ ամբողջ թվեր են: Այս դեպքում գոյություն ունի երկու ամենահավանական թիվ, և դրանց հավանականությունները հավասար են՝ $P_n(m_0) = P_n(m_0+1)$:
\end{itemize}

\textbf{Գործնական կանոն.} Ամենահավանական թիվը գտնելու համար հաճախ հարմար է հաշվել $np+p$ թիվը և վերցնել դրա ամբողջ մասը (եթե այն ամբողջ չէ)։

\end{document}